% Introduction — target 600-900 words for PLOS Comp Bio.
% Funnel: broad CPG context  →  Rybak/Danner/Zhang lineage  →
%        their stated limitations  →  our four extensions  →  paper roadmap.

\todo{4-6 paragraphs. Suggested outline below.}

\paragraph{$\bullet$ Paragraph 1 --- CPG context.}
Coordinated rhythmic limb movement in mammals is controlled primarily by
spinal cord circuitry~\cite{grahamBrown1911,kiehn2006,kiehn2016}.
Rhythm-generating (RG) circuits in the cervical and lumbar enlargements
produce flexor and extensor bursts that drive locomotor output, with the
characteristic flexor--extensor counter-phase pattern preserved across
species. \emph{The remainder of P1: a sentence on why a quantitative
computational understanding matters (e.g.\ for neural prosthetics,
recovery-after-injury models, robotics-of-locomotion).}

\paragraph{$\bullet$ Paragraph 2 --- Rybak/Danner/Zhang model lineage.}
A succession of biophysically detailed computational models has
progressively refined the architecture of the central spinal
CPG~\cite{rybak2006,shevtsova2015,danner2017,zhang2022}. These models use
Hodgkin--Huxley dynamics with persistent-sodium-driven intrinsic bursting,
typed genetic interneuron populations (V0V, V0D, V1, V2a, V3, Shox2), and
hand-curated synaptic weights. The most recent extension~\cite{zhang2022}
adds ascending long-propriospinal V3 neurons to reproduce
speed-dependent gait expression (walk, trot, gallop, bound) in mice.

\paragraph{$\bullet$ Paragraph 3 --- Acknowledged limitations.}
This lineage of models, however, makes three deliberate
simplifications, explicitly listed by~\cite{zhang2022}: (i)
\emph{``the model does not include motoneurons''}; (ii) \emph{``the
model focuses exclusively on central interactions within the spinal cord,
without considering biomechanics and the role of sensory feedback from the
limbs''}; (iii) reflex circuits mediated by Ia and Ib interneurons,
Renshaw cells, and pattern-formation interneurons are absent. In
addition, all synaptic weights in these models are \emph{hand-tuned}; no
form of synaptic plasticity is considered.

\paragraph{$\bullet$ Paragraph 4 --- Our extensions.}
Here we extend this lineage with four mechanisms:
(1) STDP-driven self-organisation of descending (BS$\to$RG) and
cutaneous (CUT$\to$RG) synapses, removing the need to hand-tune
afferent weights;
(2) closed-loop force-driven Ia spindle feedback projecting back into
the reciprocal-inhibition interneurons (InE, InF);
(3) motoneuron pools and Hill-style muscle proxies producing
graded force and length traces;
(4) externally paced sequential Ia activation modelling
heel$\to$mid$\to$toe pressure transfer during the stance phase.

\paragraph{$\bullet$ Paragraph 5 --- Scope and non-scope.}
We deliberately retain the present-paper scope to a two-leg
(left/right hindlimb) system. The four-limb fore--hindlimb
coordination captured by V3 ascending long-propriospinal neurons
in~\cite{zhang2022} is beyond the present scope. We do not attempt
to reproduce gait \emph{transitions} (walk$\to$trot$\to$gallop$\to$bound)
either: the present focus is on stability of a single paced gait under
self-organised weights, across the rat locomotor speed range
(\SIrange{600}{1400}{\milli\second} step period).

\paragraph{$\bullet$ Paragraph 6 --- Paper roadmap.}
Section~\ref{sec:methods} describes the model architecture, Izhikevich
neuron parameters, STDP rule, sensory pathways, and the paced gait
protocol. Section~\ref{sec:results} presents (i) STDP convergence and
long-term stability of the paced locomotor pattern over
\SI{120}{\second}, (ii) an ablation budget identifying which components
are necessary, and (iii) a speed sweep showing that the gait tracks
commanded step period within tens of microseconds across the rat
locomotor range. Section~\ref{sec:discussion} positions the work
relative to the Rybak lineage and discusses limitations and future
extensions.
